% Recovered antecedents. Requires amsmath, amssymb, amsthm, booktabs, longtable.
\section{Finite transport realization with phase memory}
\label{iii:hist:antecedentes}

The following developments bring together two complementary mechanisms:
a finite transfer on APP positions, modes and phase, and a collection of
kinematic and constitutive readouts on two angular channels. Each
realization retains its state space, operations and parameters. The finite
operator has the domain defined below; the enriched TPK developed in the
main text also retains its own coordinates and prolongations. Reciprocal
readouts are presented as auxiliary realizations with specific formulae
and a declared domain.

\subsection{Finite transfer on positions, modes and phase}
\label{iii:hist:transferencia}

Let $\operatorname{dr}_9(n)$ be the representative in $\{1,\ldots,9\}$
of $n\bmod9$, with $9$ representing the zero residue. For positions
$i,j\in\{0,\ldots,8\}$, consider the two tables
\[
 A_+(i,j)=\operatorname{dr}_9(i+j+2),\qquad
 A_\times(i,j)=\operatorname{dr}_9((i+1)(j+1)).
\]
The arguments evaluated here are positive integers. Consequently, the
alternative historical convention $\operatorname{dr}_9(0)=0$ does not
enter these tables. The finite state space is
\[
 \Omega=(\mathbb Z/9\mathbb Z)^2\times\{+,\times\}
 \times\{N,E,S,O\}\times\mathbb Z/3\mathbb Z,
 \qquad |\Omega|=1944.
\]
Write $x=(i,j,m,d,\tau)$, with $O$ denoting the westward direction.
The transition retains the following order: choose $d'$, move the position,
read APP in the previous mode, reduce modulo $3$, update the mode and
advance the phase. With
\[
 \delta_N=(-1,0),\quad\delta_E=(0,1),\quad
 \delta_S=(1,0),\quad\delta_O=(0,-1),
\]
the operations are
\begin{align*}
 (i',j')&=(i,j)+\delta_{d'}\pmod9,\\
 a&=A_m(i',j')\pmod3,\\
 m'&=\begin{cases}+,&a=1,\\\times,&a=2,\\m,&a=0,\end{cases}
 &\tau'&=\tau+1\pmod3.
\end{align*}
Thus $F_{d'}(x)=(i',j',m',d',\tau')$. The turning and mode-switching cost is
\[
 c_\lambda(x,d')=\mathbf1_{d'\ne d}+\lambda\mathbf1_{m'\ne m},
 \qquad\beta>0,\quad\lambda\geq0,
\]
and the operator acts on the state basis by
\begin{equation}
 T_{\beta,\lambda}e_x
 =\sum_{d'\in\{N,E,S,O\}}
 e^{-\beta c_\lambda(x,d')}e_{F_{d'}(x)}.
 \label{iii:hist:T}
\end{equation}
The parameters $\beta,\lambda$ are declared parameters of this finite
model. The reference table uses $(\beta,\lambda)=(1,1)$; the case
$\lambda=0$ belongs to the sensitivity tests.

\begin{proposition}[Equivariance of the finite antecedent]
\label{iii:hist:equivarianza}
The action
\[
 \Phi_{u,v,w}(i,j,m,d,\tau)
 =(i+3u,j+3v,m,d,\tau+w),\qquad
 (u,v,w)\in(\mathbb Z/3\mathbb Z)^3,
\]
commutes with $T_{\beta,\lambda}$.
\end{proposition}
\begin{proof}
Since $\operatorname{dr}_9(n)\equiv n\pmod3$, one has
$A_+(i,j)\equiv i+j+2$ and
$A_\times(i,j)\equiv(i+1)(j+1)\pmod3$.
Translation by multiples of three preserves both reduced readouts.
It also commutes with the chosen displacement and preserves the previous
mode, the subsequent mode and both directions. It therefore preserves
each cost indicator. Finally, adding $w$ to the phase commutes with adding
one unit. Every transition and its weight are transported to the
corresponding transition and weight; their sum in \eqref{iii:hist:T}
retains the same equality.
\end{proof}

\begin{theorem}[Complete character decomposition]
\label{iii:hist:descomposicion}
For $\omega=e^{2\pi i/3}$, the space $\mathbb C^\Omega$ decomposes
unitarily into twenty-seven subspaces of dimension $72$. Their blocks obey
\[
 T(k_a,k_b,k_t)=\omega^{k_t}T(k_a,k_b,0),
 \qquad k_a,k_b,k_t\in\{0,1,2\}.
\]
\end{theorem}
\begin{proof}
Write $i=3a+r$, $j=3b+s$, with $a,b,r,s\in\{0,1,2\}$.
Every fine fibre $u=(r,s,m,d)$ contains $27$ states $(a,b,\tau)$.
Define the isometry $U_k:\mathbb C^{72}\to\mathbb C^\Omega$ by
\[
 U_ke_u=\frac1{\sqrt{27}}
 \sum_{a,b,\tau=0}^2
 \omega^{-(k_aa+k_bb+k_t\tau)}e_{(3a+r,3b+s,m,d,\tau)}.
\]
Character orthogonality gives $U_k^*U_\ell=\delta_{k\ell}I$ and
$\sum_kU_kU_k^*=I$. If a transition crosses fine boundaries with coarse
increments $\Delta a,\Delta b$, its coefficient in $U_k^*TU_k$ is
\[
 e^{-\beta c_\lambda}\omega^{k_a\Delta a+k_b\Delta b+k_t}.
\]
This proves $TU_k=U_kT(k)$ and the operator decomposition. The temporal
increment is always one, so the final factor is common to the whole block.
Cubing removes that temporal factor, while the spatial characters retain
their role.
\end{proof}

This decomposition rests on equivariance and character orthogonality.
Nonnegativity of the trivial block also guarantees a nonnegative real
eigenvalue equal to the spectral radius. This property applies to the full
space, retaining its transient states as well.

\subsection{Spectral decomposition and parametric sensitivity}
\label{iii:hist:espectros}

Let $\rho_1\geq\rho_2$ be the two largest moduli, counted with
multiplicity, and let $g=\rho_1-\rho_2$. The following table contains
the nine spatial blocks. The preceding theorem determines all
twenty-seven temporal blocks: for every row and every $t=0,1,2$, the
dominant eigenvalue is $\omega^t\mu_1(k_a,k_b,0)$, whereas
$\rho_1,\rho_2,g$ remain unchanged. This representation explicitly
preserves all three phases. This first table fixes $\beta=\lambda=1$.

\begin{center}
\small
\begin{tabular}{ccrrrr}
\toprule
$k_a$&$k_b$&$\mu_1(k_a,k_b,0)$&$\rho_1$&$\rho_2$&$g$\\
\midrule
0&0&1.730888896&1.730888896&1.083076285&0.647812612\\
0&1&1.459962951&1.459962951&1.116843785&0.343119167\\
0&2&1.459962951&1.459962951&1.116843785&0.343119167\\
1&0&1.459962951&1.459962951&1.116843785&0.343119167\\
1&1&-1.508515719&1.508515719&1.001353790&0.507161929\\
1&2&-1.514010372&1.514010372&0.971977167&0.542033205\\
2&0&1.459962951&1.459962951&1.116843785&0.343119167\\
2&1&-1.514010372&1.514010372&0.971977167&0.542033205\\
2&2&-1.508515719&1.508515719&1.001353790&0.507161929\\
\bottomrule
\end{tabular}
\end{center}

The signs of the dominant eigenvalues and their phases are part of the
output: the four sectors with $k_a,k_b\ne0$ have a negative real leader for
$k_t=0$. Its modulus preserves the spectral rate, while its sign retains
additional information that enters the composition of iterations.

The five complete sensitivity maps, with rows $k_a$ and columns $k_b$,
are as follows; all correspond to $k_t=0$:
\begin{align*}
\rho(1,1)&=\begin{pmatrix}
1.730888896&1.459962951&1.459962951\\
1.459962951&1.508515719&1.514010372\\
1.459962951&1.514010372&1.508515719
\end{pmatrix},\\
\rho(1,0)&=\begin{pmatrix}
2.103638324&1.685504327&1.685504327\\
1.685504327&1.853245966&1.853245966\\
1.685504327&1.853245966&1.853245966
\end{pmatrix},\\
\rho(1,2)&=\begin{pmatrix}
1.656915306&1.433869155&1.433869155\\
1.433869155&1.423885161&1.432116157\\
1.433869155&1.432116157&1.423885161
\end{pmatrix},\\
\rho(0.5,1)&=\begin{pmatrix}
2.480049445&2.124801655&2.124801655\\
2.124801655&2.281861809&2.286772914\\
2.124801655&2.286772914&2.281861809
\end{pmatrix},\\
\rho(2,1)&=\begin{pmatrix}
1.226049377&1.131390710&1.131390710\\
1.131390710&1.028041508&1.028201900\\
1.131390710&1.028201900&1.028041508
\end{pmatrix}.
\end{align*}
Here $\rho(\beta,\lambda)$ denotes the map of spectral radii, not the
transfer operator. Interchanging $i$ and $j$, together with the
corresponding interchange of directions, preserves APP, transitions and
cost; these maps are therefore symmetric. Varying the parameters changes the radii while preserving the
character decomposition, which follows from symmetry of the operations.
The tables display numerical approximations to finite spectra. Transition,
equivariance and factorization identities follow from the preceding proofs;
their numerical consequences are retained together with the complete data
for all twenty-seven components.

\subsection{Reciprocal readouts and sensitivity tests}
\label{iii:hist:n6d}

Reproduction of the antecedent angular realization preserves the declared
decimal numerical measures in degrees
\[
 a_{\mathrm h}=7.297352569283802,\qquad
 c_{\mathrm h}=2.123738338968646.
\]
The angles and the corresponding dimensionless coordinate are distinguished by
\[
 A_{\mathrm h}=a_{\mathrm h}{}^\circ,\qquad
 C_{\mathrm h}=c_{\mathrm h}{}^\circ,\qquad
 \alpha_{\mathrm h}=a_{\mathrm h}/1000.
\]
These data permit a comparison with that antecedent. The constants in the
current development retain their prior APP--TRIT--TPK generation and
subsequent readout; historical reproduction does not replace that
generation by the decimals written here. The angular conversion and the
channels used are
\[
 x_{\mathrm h}=\frac{\pi a_{\mathrm h}}{180},\qquad
 y_{\mathrm h}=\frac{\pi c_{\mathrm h}}{180},\qquad
 q_\pm=e^{-(x_{\mathrm h}\pm y_{\mathrm h})},
 \qquad0<q_+<q_-<1.
\]
The first readout gives
\[
 \epsilon_{90}=q_-^{90}-q_+^{90},\qquad
 c_\pm=\frac{c_{\mathrm{ref}}}{1\pm\epsilon_{90}}.
\]
Here $c_{\mathrm{ref}}$ only normalizes the kinematic test. Reproduction
in the stated units uses the same historical reference value,
$299792458\,\mathrm{m\,s^{-1}}$.

\begin{proposition}[Harmonic mean and two-sheet product]
\label{iii:hist:lectores}
If $|\epsilon_{90}|<1$, the harmonic mean of $c_+$ and $c_-$ is
$c_{\mathrm{ref}}$. In the historical constitutive readout
\[
 \mu_{\mathrm h}^{\pm}=e^{\pm2\alpha_{\mathrm h}},\qquad
 \varepsilon_{\mathrm h}^{\pm}=e^{\mp2\alpha_{\mathrm h}},
\]
one has $\mu_{\mathrm h}^{\pm}\varepsilon_{\mathrm h}^{\pm}=1$ and
$Z_+/Z_-=e^{4\alpha_{\mathrm h}}$.
\end{proposition}
\begin{proof}
The sum $c_+^{-1}+c_-^{-1}=2/c_{\mathrm{ref}}$ removes the oriented term.
On each sheet, the two constitutive exponents sum to zero. Moreover,
$Z_\pm=\sqrt{\mu_{\mathrm h}^{\pm}/\varepsilon_{\mathrm h}^{\pm}}
=e^{\pm2\alpha_{\mathrm h}}$. The impedance ratio therefore retains
the orientation removed by the constitutive product.
\end{proof}

The second historical readout follows from convergent geometric series:
\begin{align*}
 S_a(q)&=\frac{q^a}{1-q^{3a}}=\sum_{j\geq0}q^{a(1+3j)},\\
 \delta_{a,b}(q)&=S_a(q)-S_b(q),\\
 R_{a,b}&=\frac{\delta_{a,b}(q_-)-\delta_{a,b}(q_+)}{q_-^a-q_+^a}.
\end{align*}
Convergence holds for $|q|<1$; the denominator of the last quotient is
positive for the channels and positive integers considered here. The
historical auxiliary readouts are
\[
 \zeta_{12}=1-R_{90,120}^2,\qquad
 \Delta_{4,\mathrm h}=\frac{x_{\mathrm h}-y_{\mathrm h}}{810}
 \left(1-\frac{7\pi}{12\cdot729}\right),\qquad
 \xi_{\mathrm h}=1-\frac{27}{160}\Delta_{4,\mathrm h}.
\]
The subscript $\mathrm h$ keeps this historical definition of torsion
explicit, avoiding identification with a different current readout merely
because it uses the same symbol.

Evaluation from the declared decimal strings gives
\begin{align*}
 q_+&=0.848377942576404356\ldots,&
 q_-&=0.913660151150557274\ldots,\\
 \epsilon_{90}&=0.000295169439289827959\ldots,&
 R_{90,120}&=0.933314535237713061\ldots,\\
 \zeta_{12}&=0.128923978314011665\ldots,&
 \xi_{\mathrm h}&=0.999981235497802379\ldots.
\end{align*}
The complete table of discrete perturbations is
\begin{center}
\begin{tabular}{rrr}
\toprule
$a$&$b$&$R_{a,b}$\\\midrule
90&120&0.933314535238\\
90&119&0.927013603649\\
90&121&0.939071551912\\
89&120&0.939066204789\\
91&120&0.927019444194\\
30&40&0.569215041549\\
60&80&0.834174862952\\
\bottomrule
\end{tabular}
\end{center}
Angular sensitivity distinguishes the two coordinates of the pair:
\begin{center}
\begin{tabular}{lrrr}
\toprule
Variable&Relative change&$R_{90,120}$&Change in ppm\\\midrule
$A_{\mathrm h}$&$-0.001$&0.933059250359&$-273.525022305$\\
$A_{\mathrm h}$&$+0.001$&0.933568846564&$+272.481908792$\\
$C_{\mathrm h}$&$-0.001$&0.933388163678&$+78.889202964$\\
$C_{\mathrm h}$&$+0.001$&0.933240822366&$-78.979667757$\\
\bottomrule
\end{tabular}
\end{center}

Evaluations used $80$ and $120$ digits of working precision, starting from
the declared decimal strings. Evaluation precision preserves the finite
specification of those angular data.

Finally, these readouts retain their place in the documentary genealogy.
The current constitutive operator
$V=(I-T^{90})(I-T^{120})^{-1}$ evaluates
$(1-q_\pm^{90})/(1-q_\pm^{120})$ on its two sheets. The current oriented
functional uses $12S_{90}-S_{120}$. The antecedent realization retains the readout
$S_{90}-S_{120}$ and its normalization by $q_-^{90}-q_+^{90}$.
Their relation is expressed by preserving each operation and its weights,
rather than replacing them by a single number or the same symbol.
